% =============================================================================
% CAPÍTULO 1 — Introducción General
% =============================================================================

\chapter{Introducción General}

% -----------------------------------------------------------------------------
% 1.1 CONTEXTO Y MOTIVACIÓN
% QUÉ ESCRIBIR: la importancia del mantenimiento en la industria, el costo
% del fallo no planificado, y la evolución del mantenimiento reactivo al
% predictivo. Cita a Mobley (2002) como referencia fundacional:
%   \cite{mobley2002}
% Si hablas de estrategias de mantenimiento (RCM), cita también:
%   \cite{moubray1997}
% Extensión sugerida: 2-3 párrafos.
% -----------------------------------------------------------------------------
\section{Contexto y Motivación}

% -----------------------------------------------------------------------------
% 1.2 PLANTEAMIENTO DEL PROBLEMA
% QUÉ ESCRIBIR: el mantenimiento por intervalos (preventivo) es ineficiente;
% el predictivo exige estimar la vida útil remanente (RUL) y detectar
% anomalías; se necesitan métodos precisos Y eficientes en entornos CPU-only.
% Extensión sugerida: 1-2 párrafos.
% -----------------------------------------------------------------------------
\section{Planteamiento del Problema}

% -----------------------------------------------------------------------------
% 1.3 PREGUNTA DE INVESTIGACIÓN
% QUÉ ESCRIBIR: la pregunta precisa que responde la tesis, en términos de
% viabilidad del cómputo inteligente (incluido el Predictor Gamma) para
% mantenimiento predictivo considerando precisión y costo computacional.
% Extensión sugerida: 1 párrafo.
% -----------------------------------------------------------------------------
\section{Pregunta de Investigación}

% -----------------------------------------------------------------------------
% 1.4 OBJETIVOS
% QUÉ ESCRIBIR: el objetivo general (1 párrafo) y la lista de objetivos
% específicos (enumerados con \begin{enumerate}...\end{enumerate}).
% -----------------------------------------------------------------------------
\section{Objetivos}
\subsection{Objetivo General}
\subsection{Objetivos Específicos}

% -----------------------------------------------------------------------------
% 1.5 ALCANCE Y DELIMITACIONES
% QUÉ ESCRIBIR: qué cubre la tesis (dataset C-MAPSS, 4 subconjuntos, entorno
% CPU-only, modelos de referencia + Predictor Gamma) y qué NO cubre (GPU,
% datos reales de planta, despliegue en tiempo real, optimización exhaustiva).
% Extensión sugerida: 1 párrafo.
% -----------------------------------------------------------------------------
\section{Alcance y Delimitaciones}

% -----------------------------------------------------------------------------
% 1.6 CONTRIBUCIÓN
% QUÉ ESCRIBIR: qué aporta el trabajo (evaluación multi-dimensional de
% precisión + eficiencia + recursos, y el Predictor Gamma como estimador de RUL).
% Extensión sugerida: 1 párrafo.
% -----------------------------------------------------------------------------
\section{Contribución}

% -----------------------------------------------------------------------------
% 1.7 ESTRUCTURA DEL DOCUMENTO
% QUÉ ESCRIBIR: una breve descripción de cada capítulo, del estilo:
% "El Capítulo 2 presenta el marco teórico... El Capítulo 3 describe...".
% Extensión sugerida: 1 párrafo.
% -----------------------------------------------------------------------------
\section{Estructura del Documento}
